%%=============================================================================
%% Conclusie
%%=============================================================================

\chapter{Conclusie}%
\label{ch:conclusie}

% TODO: Trek een duidelijke conclusie, in de vorm van een antwoord op de
% onderzoeksvra(a)g(en). Wat was jouw bijdrage aan het onderzoeksdomein en
% hoe biedt dit meerwaarde aan het vakgebied/doelgroep? 
% Reflecteer kritisch over het resultaat. In Engelse teksten wordt deze sectie
% ``Discussion'' genoemd. Had je deze uitkomst verwacht? Zijn er zaken die nog
% niet duidelijk zijn?
% Heeft het onderzoek geleid tot nieuwe vragen die uitnodigen tot verder 
%onderzoek?

Uit de uitgevoerde experimenten en de opbouw van het voorgestelde systeem kan geconcludeerd worden dat het mogelijk is om een functioneel anti-cheatsysteem te ontwikkelen op basis van gedragsanalyse en machine learning. In tegenstelling tot klassieke anti-cheatoplossingen, die zich voornamelijk richten op softwarematige detectie (zoals procesinjectie of geheugenmanipulatie), focust deze aanpak zich op de analyse van inputgedrag op HID-niveau.

De resultaten tonen aan dat vooral twee factoren een sterke bijdrage leveren aan de detectie van AI-gegenereerde input, namelijk de snelheid van de muisbeweging en de afgelegde afstand binnen een bepaalde tijdsinterval. Deze kenmerken vertonen bij geautomatiseerde systemen een significant ander patroon dan bij menselijke input. AI-gestuurde bewegingen zijn doorgaans consistenter, sneller en vertonen minder variatie in traject.

Hoewel deze factoren reeds een goede basis vormen voor classificatie, blijkt uit de analyse dat de complexiteit van mogelijke cheatimplementaties een belangrijke beperking vormt. Moderne AI-cheats kunnen immers bijkomende eigenschappen simuleren, zoals gebogen bewegingstrajecten, variabele versnellingen of realistische ruis in de input. Hierdoor wordt het onderscheid tussen menselijke en geautomatiseerde input minder eenduidig.

Dit betekent dat de effectiviteit van het systeem sterk afhankelijk blijft van de kwaliteit en diversiteit van de gebruikte features. Een uitbreiding met meer geavanceerde gedragskenmerken, zoals acceleratieprofielen of tijdsafhankelijke bewegingsmodellen, zou de robuustheid van het systeem verder kunnen verbeteren.

Samenvattend kan gesteld worden dat het ontwikkelde systeem aantoont dat machine learning een haalbare en veelbelovende methode is voor het detecteren van AI-gebaseerde cheats, maar dat verdere verfijning noodzakelijk is om bestand te zijn tegen meer geavanceerde en adaptieve cheattechnieken.